\documentclass[../../main.tex]{subfiles}
\begin{document}
% SOURCE p=243; printed=228; object=ch07-open-set-product-integrals; number=7.25; region=upper-middle
\begin{proposition}\label{proposition:ch07-open-set-product-integrals}
设 $\mu,\nu$ 分别为 $X,Y$ 上的 Radon 测度。若 $U$ 在 $X\times Y$ 中开，则函数 $x\mapsto\nu(U_x)$、$y\mapsto\mu(U^y)$ 分别是 $X,Y$ 上的 Borel 可测函数，且
\begin{align*}
(\mu\mathbin{\widehat\times}\nu)(U)
=\int\nu(U_x)\,\mathrm d\mu(x)=\int\mu(U^y)\,\mathrm d\nu(y).
\end{align*}
\end{proposition}
% SOURCE p=243; printed=228; object=ch07-open-set-product-integrals-proof; region=middle
\begin{proof}
令 $\mathcal F=\{f\in C_c(X\times Y):0\leqslant f\leqslant\chi_U\}$。由\cref{proposition:ch07-lsc-properties}，$\chi_U=\sup_{f\in\mathcal F}f$，所以 $\chi_{U_x}=\sup_{f\in\mathcal F}f_x$，$\chi_{U^y}=\sup_{f\in\mathcal F}f^y$。由\cref{proposition:ch07-directed-lsc-convergence}，
\begin{align*}
(\mu\mathbin{\widehat\times}\nu)(U)
&=\sup_{f\in\mathcal F}\int f\,\mathrm d(\mu\mathbin{\widehat\times}\nu),\\
\nu(U_x)&=\sup_{f\in\mathcal F}\int f_x\,\mathrm d\nu,\qquad
\mu(U^y)=\sup_{f\in\mathcal F}\int f^y\,\mathrm d\mu.
\end{align*}
由\cref{lemma:ch07-continuous-partial-integrals,proposition:ch07-lsc-properties}，两个截面测度函数下半连续，因而 Borel 可测。再次使用\cref{proposition:ch07-directed-lsc-convergence}，并结合\cref{proposition:ch07-cc-product-integrals}，得到
\begin{align*}
(\mu\mathbin{\widehat\times}\nu)(U)
&=\sup_{f\in\mathcal F}\int\left(\int f_x\,\mathrm d\nu\right)\mathrm d\mu(x)\\
&=\int\left(\sup_{f\in\mathcal F}\int f_x\,\mathrm d\nu\right)\mathrm d\mu(x)
=\int\nu(U_x)\,\mathrm d\mu(x).
\end{align*}
同理得到另一积分次序。
\end{proof}
% SUPPLEMENT: id=v14-ch07-open-set-product-integrals-details; parent=ch07-open-set-product-integrals; kind=proof-detail+motivation+hypotheses; origin=independent-derivation
\begin{remark}
\textbf{每次取上确界的条件（重构新增）。}$\mathcal F$ 含零函数；若 $f,g\in\mathcal F$，则 $\max(f,g)\in\mathcal F$，所以它向上有向。固定任一坐标后，截面族仍向上有向、非负且连续。逐点上确界等于相应的示性函数，因此对截面可应用有向下半连续收敛定理，不能以普通的\emph{序列}单调收敛定理直接替代不可数上确界。

同样，$H_f(x)=\int f_x\,\mathrm d\nu$ 非负连续，且 $H_{\max(f,g)}\geqslant\max(H_f,H_g)$，故 $\{H_f\}$ 也向上有向；其上确界正是 $\nu(U_x)$，所以再次使用该定理合法。对另一变量令 $J_f(y)=\int f^y\,\mathrm d\mu$，它们同样连续、有向且以 $\mu(U^y)$ 为上确界，故
\begin{align*}
\int\mu(U^y)\,\mathrm d\nu(y)
=\sup_{f\in\mathcal F}\int J_f\,\mathrm d\nu
=\sup_{f\in\mathcal F}\int f\,\mathrm d(\mu\mathbin{\widehat\times}\nu),
\end{align*}
给出另一等式。所有量非负，可为无穷，不涉及无穷相减。动机是利用开集示性函数的下半连续性，以连续测试从下恢复其质量；这正是开集条件的作用。这里不要求全局 $\sigma$ 有限，全部论证依赖 Radon 性和 LCH 的紧支撑逼近。
\end{remark}
% END SUPPLEMENT: id=v14-ch07-open-set-product-integrals-details
% SOURCE p=243; printed=228; object=ch07-borel-product-integrals; number=7.26; region=lower
\Needspace{9\baselineskip}
\begin{theorem}\label{theorem:ch07-borel-product-integrals}
设 $\mu,\nu$ 分别为 $X,Y$ 上的 $\sigma$ 有限 Radon 测度。若 $E\in\mathscr B_{X\times Y}$，则函数 $x\mapsto\nu(E_x)$、$y\mapsto\mu(E^y)$（由\cref{lemma:ch07-borel-sections}，它们都有定义）分别是 $X,Y$ 上的 Borel 可测函数，且
\begin{align*}
(\mu\mathbin{\widehat\times}\nu)(E)=\int\nu(E_x)\,\mathrm d\mu(x)
=\int\mu(E^y)\,\mathrm d\nu(y).
\end{align*}
此外，$\mu\mathbin{\widehat\times}\nu$ 在 $\mathscr B_X\otimes\mathscr B_Y$ 上的限制就是 $\mu\times\nu$。
\end{theorem}
% SOURCE p=243,244; printed=228-229; object=ch07-borel-product-integrals-proof; region=bottom-middle
\begin{proof}
暂时固定开集 $U\subseteq X,V\subseteq Y$，使 $\mu(U),\nu(V)$ 有限，令 $W=U\times V$。用 $\mathcal M$ 表示所有满足“$E\cap W$ 具有定理所列结论”的 $E\in\mathscr B_{X\times Y}$。则：
\begin{enumerate}[label=\roman*.,ref=(\roman*)]
\item 由\cref{proposition:ch07-open-set-product-integrals}，$\mathcal M$ 包含全部开集。
\item 若 $E,F\in\mathcal M$ 且 $F\subseteq E$，则 $E\setminus F\in\mathcal M$；特别地，$F\in\mathcal M$ 时 $F^c=(X\times Y)\setminus F\in\mathcal M$。事实上，
\begin{align*}
(\mu\mathbin{\widehat\times}\nu)(E\cap W)
&=(\mu\mathbin{\widehat\times}\nu)(F\cap W)\\
&\quad+(\mu\mathbin{\widehat\times}\nu)((E\setminus F)\cap W),
\end{align*}
截面测度也满足相应等式。对 $E\cap W,F\cap W$ 结论已成立，而所有相关测度都有限（这就是引入 $W$ 的原因），相减便得所求。
\item $\mathcal M$ 对有限不交并封闭，这来自测度的可加性。
\item $\mathcal M$ 对可数递增并封闭，因此结合第二项，也对可数递减交封闭；这是单调收敛定理的结果。
\end{enumerate}
现在令 $\mathcal E=\{A\setminus B:A,B\text{ 在 }X\times Y\text{ 中开}\}$，$\mathcal A$ 为 $\mathcal E$ 中集合的有限不交并所成族。由于
\begin{align*}
(A_1\setminus B_1)\cap(A_2\setminus B_2)&=(A_1\cap A_2)\setminus(B_1\cup B_2),\\
(A\setminus B)^c&=[(X\times Y)\setminus A]\cup[(A\cap B)\setminus\varnothing],
\end{align*}
$\mathcal E$ 是初等族，由\cref{proposition:ch01-elementary-algebra}，$\mathcal A$ 为代数。根据\cref{lemma:ch02-monotone-class}，它生成的单调类与 $\sigma$ 代数相同，即 $\mathscr B_{X\times Y}$。由前四项，并注意 $A\setminus B=A\setminus(A\cap B)$，$\mathcal M$ 包含这个单调类，因此 $\mathcal M=\mathscr B_{X\times Y}$。

由于 $\mu,\nu$ 为 $\sigma$ 有限且外正则，可以写 $X=\bigcup_{n=1}^\infty U_n$、$Y=\bigcup_{n=1}^\infty V_n$，其中 $U_n,V_n$ 开且测度有限，并可假定两序列都递增。对 $E\in\mathscr B_{X\times Y}$，前面的论证说明每个 $E\cap(U_n\times V_n)$ 满足定理结论；再用单调收敛定理，$E$ 也满足结论。

最后，若 $E\in\mathscr B_X\otimes\mathscr B_Y$，普通 Tonelli 定理给
\begin{align*}
(\mu\times\nu)(E)=\int\nu(E_x)\,\mathrm d\mu(x)
=(\mu\mathbin{\widehat\times}\nu)(E),
\end{align*}
证明完成。
\end{proof}
% SUPPLEMENT: id=v14-ch07-borel-product-integrals-details; parent=ch07-borel-product-integrals; kind=proof-detail+correction+motivation+hypotheses; origin=independent-derivation
\begin{remark}
\textbf{有限差、单调类及全空间恢复（重构新增）。}首先对开矩形用前命题有 $(\mu\mathbin{\widehat\times}\nu)(W)=\mu(U)\nu(V)<\infty$。对任意 $E$，截面 $(E\cap W)_x$ 包含于 $V$，而且在 $x\notin U$ 时为空，所以其测度至多 $\nu(V)\chi_U(x)$；另一截面至多 $\mu(U)\chi_V(y)$。这些是有限且可积的界。若 $F\subseteq E$，逐点相减便给截面差的可测性，积分的有限相减与测度的有限相减给公式。原书第二项的文字把 $F^c$ 写成 $X\setminus F$，此处环境空间是 $X\times Y$，故已明确订正为 $(X\times Y)\setminus F$。

对有限不交并，各截面仍不交，故其测度函数为有限和，积分公式相加。对 $E_n\uparrow E$，截面集合递增，测度的从下连续性使截面函数递增到目标函数；其极限 Borel，单调收敛又使积分趋于目标积分，左侧乘积测度也从下连续。递减交通过补集和递增并得到。初等族中的每个差属于 $\mathcal M$，因为两个开集 $A,A\cap B$ 都在 $\mathcal M$ 且后者包含于前者；有限不交并封闭使 $\mathcal A\subseteq\mathcal M$。因此使用单调类引理时，两项关键条件“包含代数”与“自身是单调类”都已核验。

为具体构造开覆盖，先取 $X=\bigcup_nE_n$、$\mu(E_n)<\infty$。外正则性给开集 $O_n\supseteq E_n$，$\mu(O_n)<\mu(E_n)+1<\infty$；置 $U_n=\bigcup_{j=1}^nO_j$，便得到递增有限测度开覆盖。$Y$ 同理。$W_n=U_n\times V_n\uparrow X\times Y$，所以固定 $x$ 后 $(E\cap W_n)_x\uparrow E_x$，另一个变量也一样，全部推广步骤均为单调极限，不需要减去无穷量。

动机是把开集上的积分公式扩展到所有 Borel 集。$\sigma$ 有限性与外正则性一起提供可数个有限\emph{开}矩形，而有限性使补集差的步骤合法。只知道每个截面 Borel，不保证截面测度函数 Borel；这里的单调类证明正好补足这一点。两因子 $\sigma$ 有限还保证最后一次普通 Tonelli 可用；第一个结论适用于可能更大的 $\mathscr B_{X\times Y}$，并未声称两个 $\sigma$ 代数因此相等。
\end{remark}
% END SUPPLEMENT: id=v14-ch07-borel-product-integrals-details
% SOURCE p=244; printed=229; object=ch07-radon-fubini-tonelli; number=7.27; region=lower
\begin{theorem}[Radon 乘积的 Fubini--Tonelli 定理]\label{theorem:ch07-radon-fubini-tonelli}
设 $\mu,\nu$ 分别为 $X,Y$ 上的 $\sigma$ 有限 Radon 测度，$f$ 为 $X\times Y$ 上的 Borel 可测函数。则对每个 $x,y$，$f_x,f^y$ 都是 Borel 可测函数。若 $f\geqslant0$，则 $x\mapsto\int f_x\,\mathrm d\nu$、$y\mapsto\int f^y\,\mathrm d\mu$ 分别在 $X,Y$ 上 Borel 可测。若 $f\in L^1(\mu\mathbin{\widehat\times}\nu)$，则几乎处处有 $f_x\in L^1(\nu)$、$f^y\in L^1(\mu)$，且两个部分积分函数分别属于 $L^1(\mu),L^1(\nu)$。在两种情形下，都有
\begin{align*}
\int f\,\mathrm d(\mu\mathbin{\widehat\times}\nu)
=\iint f\,\mathrm d\mu\,\mathrm d\nu
=\iint f\,\mathrm d\nu\,\mathrm d\mu.
\end{align*}
\end{theorem}
% SOURCE p=245; printed=230; object=ch07-radon-fubini-tonelli-proof; region=top
\begin{proof}
截面函数的可测性已由\cref{lemma:ch07-borel-sections}建立。其余证明与普通 Fubini--Tonelli 定理完全相同，只需以\cref{theorem:ch07-borel-product-integrals}替代\cref{theorem:ch02-product-section-integral}。
\end{proof}
% SUPPLEMENT: id=v14-ch07-radon-fubini-tonelli-details; parent=ch07-radon-fubini-tonelli; kind=omitted-proof+motivation+hypotheses; origin=independent-derivation
\begin{remark}
\textbf{从示性函数到全部函数的完整补证（重构新增）。}令 $\lambda=\mu\mathbin{\widehat\times}\nu$。对 $f=\chi_E$，前定理给出两个部分积分的 Borel 可测性及等式。对非负简单函数 $s=\sum_{j=1}^ma_j\chi_{E_j}$，有限线性组合给出同样结论。对任意非负 Borel 函数 $f$（允许取 $+\infty$），选 $0\leqslant s_n\uparrow f$ 的 Borel 简单函数序列。每个截面也有 $(s_n)_x\uparrow f_x$，所以
\begin{align*}
H_n(x)=\int(s_n)_x\,\mathrm d\nu\uparrow H(x)=\int f_x\,\mathrm d\nu.
\end{align*}
$H$ 为 Borel 函数，单调收敛给 $\int H\,\mathrm d\mu=\lim_n\int H_n\,\mathrm d\mu=\lim_n\int s_n\,\mathrm d\lambda=\int f\,\mathrm d\lambda$。交换变量，完全同理得到另一个积分等式，完成非负情形。

若 $f$ 为复值且 $\int|f|\,\mathrm d\lambda<\infty$，先对 $|f|$ 应用已证结论。函数 $A(x)=\int|f_x|\,\mathrm d\nu$ 可测，且 $\int A\,\mathrm d\mu=\int|f|\,\mathrm d\lambda<\infty$，因此 $A<\infty$ 对 $\mu$ 几乎处处成立；否则若 $\mu(A=\infty)>0$，任意常数的下界都会迫使积分无穷。故 $f_x$ 几乎处处可积。写 $f=u^+-u^-+\mathrm i(v^+-v^-)$，四个函数非负且不超过 $|f|$。分别应用非负情形后，在 $A<\infty$ 的 Borel 集上作有限差、和，得到 $H(x)=\int f_x\,\mathrm d\nu$ 的可测性；在例外零集上定义 $H=0$。由 $|H|\leqslant A$，有 $H\in L^1(\mu)$。四个有限积分的线性组合给 $\int H\,\mathrm d\mu=\int f\,\mathrm d\lambda$。对另一变量同样先用 $B(y)=\int|f^y|\,\mathrm d\mu$，在其有限集上定义部分积分、在例外零集置零，即得 $L^1(\nu)$ 性及另一等式。

动机是扩大到拓扑乘积的全部 Borel 函数。非负性允许使用无穷值而不出现不定差；复值情形的绝对可积性保证四个积分有限，才可合法相减。$\sigma$ 有限 Radon 条件用于前面的集合积分公式；对只在完备化意义下可测的代表，不能未经说明就把“每个截面 Borel”照搬，当前表述明确保留 Borel 假设。
\end{remark}
% END SUPPLEMENT: id=v14-ch07-radon-fubini-tonelli-details
\subfile{section-04-infinite-products.tex}
\end{document}
